\section{新计算范式：把 compute 推向 data}

课堂后半段回到 hardware research，讨论 in-memory、near-memory、photonics 和 compound semiconductor。它们的共同目标不是简单增加 peak FLOPS，而是减少 data movement、electrical/optical conversion 和 synchronization。读者应把这些方向看成 workload-specific design space，而不是下一代通用架构的确定名单。

\subsection{In-memory 与 near-memory computing}

\term{in-memory computing}（存内计算）让 memory device 在存储数据的位置执行部分运算；\term{near-memory computing}（近存计算）把 compute 放在 memory stack 或附近，以缩短传输路径。二者都试图缓解 von Neumann data movement，但会面对 precision、device variation、endurance、programming model 和 workload mapping 限制。

\lecturefigure{15-compute-near-memory.png}{近存/存内计算通过减少数据搬运交换能效与通用性。}{本地字幕 45:10--47:20；Sebastian et al. 2020。}

\begin{knowledgebox}{读图：为什么“少搬数据”不等于自动省一半电}
收益取决于运算是否能映射到 memory primitive、转换精度是否足够、外围 ADC/DAC 与 control 是否昂贵、数据是否需要跨层移动，以及 compiler 能否利用硬件。课堂的百分比应看作研究动机，不是对所有 workload 的保证。
\end{knowledgebox}

\subsection{术语消化：Photonic、CXL 与 collectives}

\begin{knowledgebox}{三个容易混淆的互连概念}
\begin{tabularx}{\textwidth}{L{0.19\textwidth}L{0.25\textwidth}X}
\toprule
术语 & 解决的问题 & 机制与边界 \\
\midrule
Photonic interconnect & 降低长距离高带宽传输的能耗和密度压力 & 用光传输数据，但 electro-optical conversion、laser、packaging 和 control 仍有成本。 \\
CXL & CPU、accelerator 与 memory 的 coherent attachment / pooling & 改善资源组合与共享，不等于消除远端 memory latency 或 contention。 \\
Collectives & 多 GPU 的 all-reduce、all-gather、reduce-scatter、broadcast 等集合通信 & 是 distributed training/serving 的软件语义；性能取决于 topology、message size、overlap 和 library。 \\
\bottomrule
\end{tabularx}
\end{knowledgebox}

\begin{warningbox}{Architecture religion 与 model religion 同样危险}
新架构可能在特定 kernel 上极强，却缺少 compiler、debugger、library、availability 或 multi-tenant isolation。采购和研究都应比较 end-to-end workload、migration cost 与 ecosystem maturity，而不是只依据器件演示或峰值能效。
\end{warningbox}

\teachervoice{讲者向学生发出的邀请很直接：关注 publication、prototype 与 startup signal 的交汇点。讲义把它改写为研究方法——先识别系统瓶颈，再选择可测 workload，最后确认器件收益能穿过 compiler 与 application stack。}

\subsection{本章小结}

新计算范式围绕 data movement 展开。In/near-memory、photonics、CXL 与新材料各自解决不同层次的问题；最终验收仍是端到端 workload 的能效、正确性、可编程性和可靠性。

